This chapter presents the world-model calculus in miniature.  It does not try
to prove every theorem or develop every branch of the formalization.  Instead,
it states the main semantic objects, the main exact regime, the main boundary,
and the role of typed structure and compiled inference.  A reader who studies
this chapter should be able to see how the whole book fits together before
entering the deeper technical developments.

The guiding idea is simple.  A reasoner should not treat truth-value
summaries as if they were the thing being stored or proved.  It should
carry a revisable state of knowledge, ask queries of that state, extract
answers valued in an appropriate carrier, and only then project those
answers into operational views such as strength, confidence, or interval
bounds.  This material began in the formalization of Probabilistic Logic
Networks (PLN), and along the way we found a clean, general presentation
in terms of world models.  This book develops that presentation
precisely.

\section{The Core Picture}
\label{sec:overview-core-picture}

At the highest level, the calculus has three typed parameters:
\[
  \textit{State}, \qquad \textit{Query}, \qquad V.
\]
A state is what the system currently carries.  A query is what the
system may ask.  The value type~$V$ is the carrier in which extracted
answers live.  It can be any type suitable for the application:
sufficient statistics for statistical inference, probability values,
interval-valued assessments, or even complex amplitudes for
quantum-like frameworks.

The most general world-model interface records only three operations:

\begin{lstlisting}[language={}]
class WorldModel (State Query V : Type*) where
  revise  : State -> State -> State
  empty   : State
  extract : State -> Query -> V
\end{lstlisting}

This interface is intentionally minimal.  It says only that a world model is a
revisable state space equipped with query extraction.  No additivity,
commutativity, overlap law, or probabilistic interpretation is assumed at this
top level.  Those laws belong to later subclasses and later chapters.

The book works mostly in one especially important specialization: the
\emph{additive regime}.  There, state revision is additive, the value carrier is
an additive commutative monoid, and extraction commutes with additive
revision.  This is the regime in which the profile-hom equivalence,
Bayesian-network exactness theorems, provenance-tracked revision results, and
most of the current Lean development live.  But it is not presented as the
whole theory.  It is the main exact specialization of a broader family.

\maplecourt{Maple Court Community Commons is our running example.  Each
apartment carries a private posterior state over questions such as whether a
pipe is leaking, whether the shower is running, or whether motion indicates
occupancy.  Shared spaces carry a building-level state over hallway humidity,
laundry-room availability, and elevator health.  The same architecture appears
at both scales: revise the state as reports and sensor updates arrive, query the
state for evidence, and project that evidence into views only when a decision
must be made.}

\section{Revision Regimes and the Additive Powerhouse}
\label{sec:overview-regimes}

The calculus does not assume that all revision behaves the same way.  It is
useful to distinguish several regimes.

\begin{description}
\item[Additive revision.]  Independent updates compose by addition.
  In Maple Court: two independent sensor reports about a pipe leak
  add their evidence counts.

\item[Overlap-aware revision.]  When two sources share information,
  naive addition double-counts.  In Maple Court: two apartments'
  humidity sensors both pick up the same hallway vent.

\item[Trust-gated revision.]  The merge depends on source authority.
  In Maple Court: a tenant's verbal report contradicts an automated
  sensor; the system must weigh trust.

\item[Coalition revision.]  Not every merge is valid; consensus
  requires an actionable coalition.  In Maple Court: a majority of
  apartments must agree before updating the building's elevator-health
  assessment.
\end{description}

Only the first of these is developed theoremically in full generality in the
present book.  The others already appear in the typed family and in the
formalization landscape, but they are treated as frontier directions rather than
as fully developed kernels.

The additive regime is captured by the following shape:

\begin{lstlisting}[language={}]
class AdditiveWorldModel (State Query V : Type*)
    [AddCommMonoid State] [AddCommMonoid V]
    extends WorldModel State Query V where
  revise_eq_add : forall W1 W2, revise W1 W2 = W1 + W2
  empty_eq_zero : empty = 0
  extract_add   : forall W1 W2 q,
    extract (W1 + W2) q = extract W1 q + extract W2 q
  extract_zero  : forall q, extract 0 q = 0
\end{lstlisting}

The decisive law is \code{extract\_add}.  It says that in the additive regime,
revising states and then extracting is the same as extracting first and then
combining answers in the carrier~$V$.  This is the exact algebraic condition
behind the cleanest part of the calculus.

\begin{theorem}[Additive answer-profile form]
\label{thm:overview-profile-hom}
For additive world models, currying the extractor yields an additive monoid
homomorphism
\[
  \widehat{\mathrm{extract}} : \textit{State} \to^+ (\textit{Query} \to V),
  \qquad
  W \mapsto (q \mapsto \mathrm{extract}(W,q)).
\]
Conversely, every additive monoid homomorphism
\(\textit{State} \to^+ (\textit{Query} \to V)\) determines an additive world
model.  Thus additive world models are exactly answer-profile generators.
\end{theorem}

This theorem is the categorical and algebraic center of the additive regime.
It says that, once additivity is assumed, a world model is one
structure-preserving arrow into an answer-profile object.  The extensional
model \(\textit{Query} \to V\) is terminal among additive world models over the
same query space: every additive world model maps canonically to its answer
profile.

This does \emph{not} trivialize the subject.  The hard work is in choosing
expressive state spaces, proving when compiled shortcuts are exact, and
showing where local summaries fail.  But it does make the kernel visible.

\section{Evidence Carriers}
\label{sec:overview-carriers}

Evidence is not one fixed datatype.  The general pattern is that a carrier
stores sufficient statistics for a conjugate family or other compositional value
system, and revision in the additive regime is addition of those sufficient
statistics.

Three carriers are central.

\begin{description}
\item[Binary evidence.]  A two-channel carrier recording positive and negative
  support.  This is the historical PLN carrier, now treated as one instance
  among others.
\item[Dirichlet evidence.]  A multi-valued carrier recording category counts.
  This is the right instance for questions with more than two alternatives.
\item[Normal--Gamma evidence.]  A continuous-domain carrier storing the
  sufficient statistics needed for unknown mean and variance.
\end{description}

In Maple Court these arise naturally.  Leak alarms and motion detections are
binary.  Elevator health and laundry status are multi-valued.  Hallway
humidity and temperature are continuous.  The calculus does not change across
these cases; only the carrier does.

The additive regime explains why revision is so simple in these families:
counts add, sufficient statistics add, and the extractor commutes with that
addition.  Some additional algebraic structure can then be studied on top of
particular carriers.  For binary evidence, the current development includes a
quantale tensor and a Heyting structure.  Those are mathematically real, but
they are \emph{binary-specific} and do not generalize to all carriers.

\section{States Are Primary: Revise, Query, Forget, Explain}
\label{sec:overview-four-verbs}

The book revolves around four verbs:
\emph{revise}, \emph{query}, \emph{forget}, and \emph{explain}.

Revision and querying are the minimal kernel operations.  Forgetting and
explanation show why the kernel stores states rather than floating truth values.
A state can be updated, partially retracted, audited, replayed, and explained in
a way that a naked scalar view cannot.

\maplecourt{Suppose three hallway humidity alerts arrive, then one sensor is
found to be faulty.  If the system carries only a strength estimate such as
$0.83$, it is unclear how to retract just the faulty source.  If the system carries
a revisable state with provenance, forgetting can remove the appropriate
contribution while preserving answers outside the forgotten scope.  The same
state can also explain why the building robot now believes the hallway needs
mopping: not because the confidence number crossed an arbitrary threshold in
the abstract, but because specific reports and sensors revised the building
state in a tracked way.}

This is why provenance is not ``mere metadata.''  It is part of the structure of
stateful reasoning.  Likewise, forgetting is not an optional convenience.  It is a
semantic operation that clarifies what kind of proof object a world-model
state really is.

\section{Information, Support, and Views}
\label{sec:overview-views}

One of the deepest distinctions in the book is that evidence can improve in
more than one sense.  It can become \emph{more informed}, and it can become
\emph{more supportive} of a query.  These are not the same order.

For binary evidence, compare
\[
  e_1 = (1,0),
  \qquad
  e_2 = (1,1).
\]
The second state is more informative than the first in the ordinary
information-order sense: it contains at least as much positive and at least as
much negative evidence.  But it is less supportive of the positive claim,
because the negative channel has increased as well.  A strength view therefore
decreases even as total evidence increases.

This distinction drives the architecture.

\begin{proposition}[Views are not additive semantics]
\label{prop:overview-views-not-semantic-core}
In the current development, total evidence is additive, but central operational
views such as strength are not.  More generally, useful readouts from evidence
may be informative, calibrated, or decision-relevant without preserving the
revision law of the underlying carrier.
\end{proposition}

So a view is not a proof object.  It is a projection.  This is why the book keeps
views after evidence and state operations, not before them.  The semantics lives
in state and extraction; views are derived interfaces for display, thresholding,
and action.

\section{Natively Typed World Models}
\label{sec:overview-typed}

The calculus is typed from the beginning.  Its primary sorts are already
visible in the kernel:
\[
  \textit{State}, \qquad \textit{Query}, \qquad V.
\]
What later chapters develop in more detail is not a second semantics but a
full native typed realization of the same calculus: a rewrite family whose
terms, rules, and native types expose the same mathematical structure from
inside.

In the additive regime, the core rewrite system is organized by a 2\(\times\)2
calculus square: raw versus guarded, and plain versus congruence-aware.
That square gives the current formalized typed face of the additive calculus.
Later, the typed development expands this to the 13-axis vertex family,
OSLF modal operators, intrinsic adequacy, and native type synthesis.
The encoding is proved biconditionally adequate: the typed and algebraic
presentations agree exactly.  For any type~$T$, multiset ensembles
give a canonical world model---suggesting the framework may be
universal for typed rewrite systems.

Logical structure appears wherever the query space carries it.  Quantified and
higher-order queries are not presented as ``extensions'' to a lower theory.
They are simply instances of the general \textit{Query} parameter when one wants
logical structure rich enough to express those questions.

\section{Compiled Inference and the PLN Connection}
\label{sec:overview-compiled}

The calculus does not treat deduction, induction, and abduction as primitive
truth-value algebra.  It treats them as \emph{compiled rewrites}: certified
shortcuts that are exact under explicit side conditions within declared model
classes.

This is where the book most visibly overlaps with the PLN tradition.  The
derived link calculus and the compiled Bayesian-network rules formalize,
sharpen, and delimit the classical PLN rule family.  The point is not to deny
that those rules can be useful.  The point is to state exactly when they are
correct, what they compute, and what they leave out.

\maplecourt{Suppose the building model contains the chain
\(\textit{PipeLeak} \to \textit{WallHumidity} \to \textit{MoldRisk}\).
A maintenance agent wants to answer a leak-related query quickly.  Rather
than recomputing the full posterior from scratch, the system may use a
compiled chain rule justified by the BN factorization and screening-off
assumptions attached to that model class.  When those side conditions hold,
the compiled answer agrees with the full extractor on the covered query
family.  When they do not hold, the rewrite is not licensed.}

That is the modernized meaning of a PLN rule: not a free-standing law
of truth values, but a typed and side-conditioned shortcut to a result
already determined by the underlying world model.

The most exact form of chaining propagates the \emph{full posterior
distribution} through the chain---not just scalar strength and
confidence.  This distributional approach is the primary one; scalar
chaining (propagating only summary statistics) is a useful but lossy
approximation whose cost is made precise by five chaining-limitation
theorems and a distributional dominance result
(Chapter~\ref{ch:error-transport}).

\section{The Boundary: Why Local Chaining Cannot Be Complete}
\label{sec:overview-boundary}

The book's main negative result explains why compiled inference must be
handled so carefully.

\begin{theorem}[No-go for unrestricted local link completeness]
\label{thm:overview-no-go}
In general, complete link evidence cannot be recovered from local per-link
surrogates alone.  Two world-model states can agree on all chosen local link
summaries while differing on the complete evidence needed for the target
query.
\end{theorem}

This boundary theorem is not an embarrassment.  It is a design theorem.  It
says that unrestricted local chaining cannot replace a complete semantics.  It
also explains why structural side conditions matter so much: compiled shards
become exact only inside model classes that carry enough global structure for
those shortcuts to commute with full extraction.

The no-go theorem is also one reason the non-additive frontier matters.  When
sources overlap, correlations matter, or trust gates revision, naive local
summary propagation is even less likely to preserve the required structure.
The additive regime is powerful precisely because it identifies a large,
tractable class in which the profile-hom story, carrier algebra, and compiled
rules all fit cleanly together.

\section{A Compact End-to-End Picture}
\label{sec:overview-end-to-end}

The overview can now be summarized in one pipeline:

%% Safe tikz width: ~148mm text width for 11pt report.
%% This diagram is ~155mm — slightly wider, centered symmetrically in margins.
\medskip
\noindent
\makebox[\textwidth][c]{%
\begin{tikzpicture}[
  >=Latex,
  every node/.style={font=\footnotesize},
  box/.style={draw, rounded corners=3pt, minimum width=30mm,
              minimum height=9mm, align=center}
]
  \node[box, fill=blue!8] (state) {state $W$};
  \node[box, fill=green!8, right=10mm of state] (query) {query $q$};
  \node[box, fill=yellow!12, right=10mm of query] (extract) {extract$(W,q)$};
  \node[box, fill=orange!12, right=10mm of extract] (view) {view / decision};
  \node[box, fill=purple!8, below=10mm of extract] (rewrite) {compiled rewrite};

  \draw[->, thick] (state) -- (query);
  \draw[->, thick] (query) -- (extract);
  \draw[->, thick] (extract) -- (view);
  \draw[->, thick] (state) |- (rewrite);
  \draw[->, thick] (rewrite) -- node[right, font=\scriptsize]
    {$\Sigma$-shortcut} (extract);
\end{tikzpicture}}
\medskip

The state is primary.  Queries interrogate the state.  Extraction returns values
in a carrier.  Views project from those values.  Compiled rewrites accelerate
families of queries under explicit side conditions.  Boundaries tell us where
such shortcuts cannot be complete.

That is the whole book in compressed form.

\section{How the Rest of the Book Unfolds}
\label{sec:overview-rest}

The remaining chapters slow this compact picture down.

\begin{itemize}
\item The next chapters develop the calculus and its additive regime in full
  detail, prove the profile-hom equivalence, and state the precise typed rules.
\item The carrier chapter develops the generic evidence pattern and the
  formalized instances.
\item The state-operations chapter develops forgetting, provenance,
  explanation, and conservation.
\item The views chapter explains why information order and support order
  must be distinguished, and why views are not proof objects.
\item The typed-development chapter develops the 13-axis family, native
  types, OSLF bridge, and the GSLT horizon.
\item The compiled-inference chapters show how the PLN rule family and the
  BN exactness results become side-conditioned rewrites.
\item The final parts map the current scope, validate the framework on
  substantial benchmarks, and mark the frontier: non-additive regimes,
  measure-theoretic grounding, higher-order completeness, GSLT universality,
  and the probability hypercube.
\end{itemize}

The intended reading experience is therefore two-layered.  This overview
shows the crystal.  The rest of the book turns the crystal in the light,
proves its faces, and follows each face into deeper mathematics and larger
applications.

\paragraph{PLN heritage.}
This book is not a historical commentary on PLN, but neither is it detached
from that lineage.  Its derived link calculus and compiled Bayesian-network
rules formalize and extend the classical probabilistic deduction,
induction, abduction, and revision story of PLN.  Where that earlier rule
family is correct, the present calculus proves it correct under explicit
assumptions.  Where it is incomplete or semantically unstable, the present
calculus states the boundary and replaces the heuristic with a typed,
side-conditioned construction.
